% Generated semantic source for AI Almanach v3.
% Edit v3 directly after initial generation; do not overwrite
% editorial changes by rerunning the converter casually.

% ==== Reinforcement Learning =============================================================================
    \chapter{Reinforcement Learning}

\chapterlead{How should an agent learn when its actions change later observations?}{Model}{Act}{Update}{Evaluate}{reinforcement-learning}

    % -------------------------------------------------------
\label{ch:reinforcement-learning}

\begin{learningobjectives}
After working through this chapter you should be able to:
\begin{itemize}
  \item compute a discounted return and convert a discount factor into an
        effective planning horizon;
  \item propagate values backwards through a small MDP by hand;
  \item execute one Q-learning update and say what each term contributed;
  \item explain why SARSA and Q-learning behave differently near a cliff;
  \item compare $\epsilon$-greedy and UCB numerically on the same arm counts;
  \item state why a constant $\epsilon$ never stops paying for exploration;
  \item distinguish a correlation, an intervention, and a counterfactual;
  \item detect Simpson's paradox in a real table of outcomes.
\end{itemize}
\end{learningobjectives}

\begin{plainlanguage}
Everything before this chapter learned from answers that were handed over.
Reinforcement learning has no answers --- only consequences, and usually late
ones. An agent acts, the world changes, and eventually a number arrives saying
roughly how things are going. The whole subject is the problem of assigning
credit: which of the last two hundred decisions earned that number? Everything
here --- value functions, discounting, temporal-difference updates --- is
machinery for answering that one question.
\end{plainlanguage}

\begin{engineernote}
RL is the most oversold technique in this book and the one most likely to waste
a year. It needs an enormous number of interactions, a reward you can actually
compute, and either a simulator or a very tolerant environment. Before choosing
RL, ask whether the problem is really sequential: if each decision does not
change the situation the next decision faces, you have a supervised or bandit
problem, and both are dramatically cheaper to solve and to verify.
\end{engineernote}

    \section{Introduction to Reinforcement Learning}


\begin{v3topic}{What is Reinforcement Learning?}
        Reinforcement Learning (RL) is a learning paradigm in which an
        \textbf{agent} interacts with an \textbf{environment} to maximise
        cumulative \textbf{reward}. Unlike supervised learning (labelled
        examples) or unsupervised learning (structure discovery), the agent
        receives only a scalar reward signal --- often delayed --- and must learn
        by \emph{trial and error}.\textsuperscript{\cite[][p.~323]{hurbansGrokkingArtificialIntelligence2020}}

        The core challenge is the \textbf{exploration--exploitation trade-off}:
        should the agent exploit its best-known strategy or explore to
        discover potentially better ones?

\end{v3topic}
\begin{v3topic}{The Reward Hypothesis}
        \textbf{Reward hypothesis}: all goals can be expressed as the
        maximisation of a scalar cumulative reward signal.

        \begin{itemize}
            \item Reward encodes the \emph{objective}, not the \emph{strategy}
            \item Reward may be immediate or heavily delayed
            \item The agent is not told \emph{how} to act --- only \emph{how well}
        \end{itemize}
        \textsuperscript{\cite[][p.~324]{hurbansGrokkingArtificialIntelligence2020}}

\end{v3topic}
\begin{v3topic}{Agent--Environment Interaction Loop}
        At each discrete timestep $t$ the agent receives state $s_t$, selects
        action $a_t$, and the environment returns reward $r_{t+1}$ and next
        state $s_{t+1}$. This produces a \textbf{trajectory}
        $\tau = (s_0, a_0, r_1, s_1, a_1, r_2, \ldots)$.%
        \textsuperscript{\cite[][p.~325]{hurbansGrokkingArtificialIntelligence2020}}
\visualseparator
        \begin{center}
        \begin{tikzpicture}[
            box/.style={draw, rounded corners=4pt, minimum width=2.2cm,
                        minimum height=0.9cm, align=center, font=\small},
            arr/.style={-Stealth, thick}
        ]
            \node[box, fill=LimeGreen!25] (agent) {Agent};
            \node[box, fill=orange!20, right=4.5cm of agent] (env) {Environment};
            \draw[arr] (agent.north) to[bend left=30]
                node[above, font=\footnotesize]{action $a_t$} (env.north);
            \draw[arr] (env.south) to[bend left=30]
                node[below, font=\footnotesize]{state $s_{t+1}$, reward $r_{t+1}$}
                (agent.south);
        \end{tikzpicture}
\captionof{figure}{The agent--environment loop. The agent never receives a
correct action, only a reward that may arrive many steps after the decision
that earned it --- which is the credit-assignment problem the rest of this
chapter exists to solve.}
\label{fig:rl-agent-environment-loop}
        \end{center}

\end{v3topic}
\begin{v3topic}{Core Components}
        \begin{description}[leftmargin=0.5cm, labelindent=0pt]
            \item[\textbf{Agent}] The learner and decision-maker
            \item[\textbf{Environment}] Everything external to the agent
            \item[\textbf{State $s$}] Representation of the current situation
            \item[\textbf{Action $a$}] A choice available to the agent
            \item[\textbf{Reward $r$}] Scalar feedback from the environment
            \item[\textbf{Policy $\pi$}] Mapping: state $\to$ action (or distribution)
            \item[\textbf{Value $V(s)$}] Expected cumulative future reward from $s$
            \item[\textbf{Model}] Internal model of environment dynamics (optional)
        \end{description}
        \textsuperscript{\cite[][p.~326]{hurbansGrokkingArtificialIntelligence2020}}

\end{v3topic}
\begin{v3topic}{RL vs.\ Supervised vs.\ Unsupervised}
        \begin{tblr}{width=\linewidth,
            colspec={X[l]X[l]X[l]},
            hlines, vlines,
            row{1}={font=\bfseries\small, bg=LimeGreen!20},
            row{2-Z}={font=\small},
        }
            Supervised & Unsupervised & Reinforcement \\
            Labelled data & Unlabelled data & Interactions \\
            Immediate answer & No feedback & Delayed reward \\
            Generalise & Discover structure & Maximise return \\
            No exploration & No exploration & Exploration essential \\
        \end{tblr}
\captionof{table}{The three learning paradigms compared by what supervision each receives. Reinforcement learning's distinguishing feature is that feedback is evaluative and delayed rather than instructive and immediate.}
\label{tab:rl-paradigms}
        \textsuperscript{\cite[][p.~323]{hurbansGrokkingArtificialIntelligence2020}}

\end{v3topic}


    % -------------------------------------------------------
    \section{Markov Decision Processes}


\begin{v3topic}{The Markov Property}
        A state $s_t$ satisfies the \textbf{Markov property} if:
        \begin{equation}
            P(s_{t+1} \mid s_t, a_t) = P(s_{t+1} \mid s_0, a_0, \ldots, s_t, a_t)
        \end{equation}
        The future depends only on the \emph{current} state and action, not
        on the full history --- the state is a \emph{sufficient statistic} of
        the past. This memoryless property enables tractable dynamic
        programming solutions.\textsuperscript{\cite[][p.~140]{barberBayesianReasoningMachine2012}}

\end{v3topic}
\begin{v3topic}{MDP: Formal Definition}
        A \textbf{Markov Decision Process} is the tuple
        $\mathcal{M} = (\mathcal{S}, \mathcal{A}, P, R, \gamma)$:
        \begin{itemize}
            \item $\mathcal{S}$: finite state space
            \item $\mathcal{A}$: finite action space
            \item $P(s' \mid s, a)$: transition probability
            \item $R(s, a)$: expected immediate reward
            \item $\gamma \in [0,1]$: discount factor
        \end{itemize}
        The agent's goal is to find policy $\pi^*$ that maximises expected
        discounted return from every starting
        state.\textsuperscript{\cite[][p.~141]{barberBayesianReasoningMachine2012}}

\end{v3topic}
\begin{v3topic}{Cumulative Discounted Return $G_t$}
        The \textbf{return} at time $t$ is the discounted sum of all future
        rewards:\textsuperscript{\cite[][p.~142]{barberBayesianReasoningMachine2012}}
\visualseparator
        \begin{equation}
            G_t = \sum_{k=0}^{\infty} \gamma^k\, r_{t+k+1}
                = r_{t+1} + \gamma r_{t+2} + \gamma^2 r_{t+3} + \cdots
        \end{equation}
        Discount factor $\gamma$ controls the time-horizon:
        $\gamma = 0$ (myopic, only immediate reward);
        $\gamma \to 1$ (far-sighted, all future equally weighted).
        For $\gamma < 1$, the sum converges even for infinite horizons.

\begin{workedexample}[What the discount factor actually chooses]
\emph{A finite return.} Rewards $r_1 = 1$, $r_2 = 2$, $r_3 = 3$ with
$\gamma = 0.5$:
\[
  G = 1 + 0.5(2) + 0.5^2(3) = 1 + 1 + 0.75 = 2.75 .
\]

\emph{An infinite one.} A constant reward of $1$ forever:
\[
  G = \sum_{k=0}^{\infty}\gamma^k = \frac{1}{1-\gamma}
  \quad\Rightarrow\quad
  \gamma = 0.9 \Rightarrow 10,
  \qquad \gamma = 0.99 \Rightarrow 100 .
\]

\emph{The interpretation that matters.} $1/(1-\gamma)$ is the \keyterm{effective
horizon} --- roughly how many steps ahead the agent cares about. At
$\gamma = 0.9$ the agent is planning about ten steps out; at $\gamma = 0.99$,
about a hundred. So $\gamma$ is not a numerical convenience, it is a statement
about how far-sighted you want the behaviour to be, and changing it from
$0.9$ to $0.99$ changes the problem, not just the arithmetic.
\end{workedexample}

\end{v3topic}
\begin{v3topic}{Value Functions}
        \textbf{State-value function} under policy $\pi$:
        \begin{equation}
            V^{\pi}(s) = \mathbb{E}_{\pi}\!\left[G_t \mid S_t = s\right]
        \end{equation}
        \textbf{Action-value (Q) function} under policy $\pi$:
        \begin{equation}
            Q^{\pi}(s, a) = \mathbb{E}_{\pi}\!\left[G_t \mid S_t = s, A_t = a\right]
        \end{equation}
        $V^{\pi}(s)$ is the expected return following $\pi$ from $s$;
        $Q^{\pi}(s,a)$ additionally fixes the first action.
        Optimal: $V^*(s) = \max_\pi V^\pi(s)$.\textsuperscript{\cite[][p.~143]{barberBayesianReasoningMachine2012}}

\end{v3topic}
\begin{v3topic}{Bellman Equations}
        \textbf{Bellman expectation} decomposes the value function
        recursively:\textsuperscript{\cite[][p.~145]{barberBayesianReasoningMachine2012}}
        \begin{equation}
            V^{\pi}(s) = \sum_a \pi(a|s) \sum_{s'} P(s'|s,a)
                         \bigl[R(s,a) + \gamma V^{\pi}(s')\bigr]
        \end{equation}
        \textbf{Bellman optimality} characterises the optimal value:
        \begin{equation}
            V^*(s) = \max_a \sum_{s'} P(s'|s,a)
                     \bigl[R(s,a) + \gamma V^*(s')\bigr]
            \label{eq:rl-bellman-optimality}
        \end{equation}

\begin{plainlanguage}
The Bellman equation says something almost obvious once you hear it: the value
of where you are equals the reward you get for your next step plus the
discounted value of wherever that step lands you. It turns one impossible
question --- ``what is this whole future worth?'' --- into a one-step question
asked over and over.
\end{plainlanguage}

\begin{workedexample}[Value propagating backwards along a chain]
Four states in a line. From each state the only action moves right. All rewards
are $0$ except entering state $4$ (terminal), which pays $+1$. Let
$\gamma = 0.9$.

Work backwards from the end, because that is the only place a known value
exists:
\[
\begin{aligned}
  V(4) &= 0 &&\text{(terminal)},\\
  V(3) &= 1 + 0.9\,V(4) = 1 + 0 &&= 1.000,\\
  V(2) &= 0 + 0.9\,V(3) = 0.9(1) &&= 0.900,\\
  V(1) &= 0 + 0.9\,V(2) = 0.9(0.9) &&= 0.810 .
\end{aligned}
\]

\emph{Read the numbers.} No state except the last one ever pays anything, yet
every state has a positive value --- the reward has propagated backwards,
decaying by a factor $\gamma$ per step of distance. That decaying trail is what
lets an agent in state $1$ prefer moving right without ever having reached the
goal in this calculation. It is also why sparse rewards are hard: at
$\gamma = 0.9$ and a goal $100$ steps away, the value at the start is
$0.9^{99} \approx 3\times10^{-5}$ --- numerically indistinguishable from
nothing.
\end{workedexample}

\begin{figure}[htbp]
\centering
\begin{tikzpicture}[
  st/.style={circle,draw=AlmanachTeal,fill=AlmanachMist,minimum size=11mm,
    font=\sffamily\small},
  term/.style={st,fill=AlmanachOchre!30,draw=AlmanachOchre},
  arr/.style={-{Stealth[length=1.8mm]},thick,draw=AlmanachInk!55},
  back/.style={-{Stealth[length=1.8mm]},thick,draw=AlmanachRed,dashed},
  lbl/.style={font=\sffamily\scriptsize,color=AlmanachInk!75},
  val/.style={font=\sffamily\scriptsize,color=AlmanachTeal!75!black}]
  \foreach \i/\n in {1/1,2/2,3/3}{
    \node[st] (s\n) at (\i*24mm-24mm,0) {$s_\n$};}
  \node[term] (s4) at (72mm,0) {$s_4$};
  \foreach \a/\b in {s1/s2,s2/s3}{\draw[arr] (\a) -- (\b) node[midway,above,lbl] {$r{=}0$};}
  \draw[arr] (s3) -- (s4) node[midway,above,lbl,color=AlmanachOchre] {$r{=}{+}1$};
  \node[val,below=3mm of s1] {$0.810$};
  \node[val,below=3mm of s2] {$0.900$};
  \node[val,below=3mm of s3] {$1.000$};
  \node[val,below=3mm of s4] {$0$};
  \draw[back] ([yshift=8mm]s4.north) -- ([yshift=8mm]s1.north)
    node[midway,above,lbl,color=AlmanachRed]
    {values are computed in this direction};
\end{tikzpicture}
\caption{Value propagating backwards from the only rewarded transition. The
agent moves left to right; the \emph{information} travels right to left,
shrinking by the discount factor at each step. Sparse-reward problems are hard
precisely because this trail becomes numerically invisible over long
distances.}
\label{fig:rl-value-propagation}
\end{figure}

\end{v3topic}
\begin{v3topic}{Value Iteration}
        A \textbf{dynamic programming} algorithm for finding $V^*$ when the
        model $(P, R)$ is
        known:\textsuperscript{\cite[][p.~153]{barberBayesianReasoningMachine2012}}
        \begin{enumerate}
            \item Initialise $V(s) \leftarrow 0\ \forall s$
            \item Repeat until $\|V_{k+1} - V_k\|_\infty < \epsilon$:
            \begin{equation}
                V_{k+1}(s) \leftarrow \max_a \sum_{s'} P(s'|s,a)
                           \bigl[R(s,a) + \gamma V_k(s')\bigr]
            \end{equation}
            \item Extract $\pi^*(s) = \arg\max_a \sum_{s'} P(s'|s,a)[R + \gamma V^*]$
        \end{enumerate}
        Convergence guaranteed (Bellman operator is a contraction).

\end{v3topic}
\begin{v3topic}{MDP State Transitions}
        MDP transitions depend on both state and action. Below: three states
        with two actions ($a_1, a_2$), illustrating stochastic transitions.
\visualseparator
        \begin{center}
        \begin{tikzpicture}[
            state/.style={circle, draw, minimum size=0.9cm, font=\small},
            arr/.style={-Stealth, thick},
            every node/.style={font=\small}
        ]
            \node[state, fill=LimeGreen!20] (s1) {$s_1$};
            \node[state, fill=orange!20, right=2.5cm of s1] (s2) {$s_2$};
            \node[state, fill=blue!15, below=1.5cm of s2] (s3) {$s_3$};
            \draw[arr, bend left=15] (s1) to
                node[above, font=\footnotesize]{$a_1,\,p{=}0.8$} (s2);
            \draw[arr] (s1) to
                node[left, font=\footnotesize]{$a_1,\,p{=}0.2$} (s3);
            \draw[arr] (s1) to[loop left]
                node[left, font=\footnotesize]{$a_2$} (s1);
            \draw[arr] (s2) to
                node[right, font=\footnotesize]{$a_1$} (s3);
            \draw[arr, bend right=20] (s3) to
                node[below, font=\footnotesize]{$a_2$} (s1);
        \end{tikzpicture}
\captionof{figure}{Transitions in a Markov decision process. Because the next
state depends only on the current state and action, the entire history can be
discarded --- which is what makes dynamic programming possible and what makes
choosing the state representation the critical modelling decision.}
\label{fig:rl-mdp-transitions}
        \end{center}

\end{v3topic}
\begin{v3topic}{Policy Types}
        \begin{description}[leftmargin=0.5cm]
            \item[\textbf{Deterministic}] $\pi : \mathcal{S} \to \mathcal{A}$
                --- one action per state
            \item[\textbf{Stochastic}] $\pi(a \mid s) = P(A{=}a \mid S{=}s)$
                --- probability distribution over actions
            \item[\textbf{Optimal $\pi^*$}] maximises $V^\pi(s)$ for every $s$;
                always exists as deterministic in finite MDPs
            \item[\textbf{Greedy}] $\pi(s) = \arg\max_a Q(s,a)$
                --- exploits current value estimate
            \item[\textbf{$\epsilon$-greedy}] greedy with probability $1{-}\epsilon$,
                random with probability $\epsilon$
        \end{description}
        \textsuperscript{\cite[][p.~144]{barberBayesianReasoningMachine2012}}

\end{v3topic}
\begin{v3topic}{Markov Chain Monte Carlo (MCMC)}
        \textbf{MCMC} uses Markov chains to draw samples from intractable
        distributions $p(x)$.

        \textbf{Metropolis-Hastings}:
        \begin{enumerate}
            \item Propose $x' \sim q(x' \mid x)$
            \item Compute acceptance ratio:
                  $\alpha = \min\!\left(1,\,
                  \frac{p(x')\,q(x \mid x')}{p(x)\,q(x' \mid x)}\right)$
            \item Accept $x'$ with probability $\alpha$; else stay at $x$
        \end{enumerate}
        At stationarity, samples follow $p(x)$.
        Applications: Bayesian posterior inference, policy
        evaluation.\textsuperscript{\cite[][p.~160]{barberBayesianReasoningMachine2012}}

\end{v3topic}


    % -------------------------------------------------------
    \section{Bandit Problems}


\begin{v3topic}{The Multi-Armed Bandit}
        The \textbf{multi-armed bandit} (MAB) is an RL setting with
        \emph{no state transitions}: $k$ actions (arms), each yielding reward
        from an unknown distribution with mean $\mu_a$.

        Goal: minimise cumulative \textbf{regret}:
        \begin{equation}
            R_T = T\,\mu^* - \sum_{t=1}^{T} r_t
        \end{equation}
        where $\mu^* = \max_a \mu_a$.
        The fundamental tension is \textbf{exploration vs.\ exploitation}:
        exploit the best-known arm or explore to reduce
        uncertainty.\textsuperscript{\cite[][p.~335]{hurbansGrokkingArtificialIntelligence2020}}

\end{v3topic}
\begin{v3topic}{$\epsilon$-Greedy Strategy}
        With probability $\epsilon$ choose a random arm (explore); with
        probability $1 - \epsilon$ choose $\arg\max_a Q(a)$ (exploit).

        Incremental mean update after observing reward $r$:
        \begin{equation}
            Q(a) \leftarrow Q(a) + \frac{1}{n}\bigl[r - Q(a)\bigr]
        \end{equation}
        \begin{itemize}
            \item Simple and robust; $\epsilon$ often decayed over time
            \item Wastes exploration uniformly across all arms, even
                  well-understood ones
        \end{itemize}
        \textsuperscript{\cite[][p.~336]{hurbansGrokkingArtificialIntelligence2020}}

\begin{workedexample}[What a constant $\epsilon$ costs, forever]
Two arms with true means $\mu_1 = 0.6$ (the best) and $\mu_2 = 0.2$. Suppose
the agent has already learned which is which, and runs $\epsilon$-greedy with
$\epsilon = 0.1$.

\emph{Action probabilities.} The greedy arm is chosen either by exploitation
($1-\epsilon = 0.9$) or by landing on it during random exploration
($\epsilon/2 = 0.05$), so $P(\text{arm 1}) = 0.95$ and
$P(\text{arm 2}) = 0.05$.

\emph{Expected reward per pull.}
\[
  0.95(0.6) + 0.05(0.2) = 0.570 + 0.010 = 0.580 .
\]

\emph{Regret.} The best possible is $0.6$, so the agent loses $0.020$ per pull
--- $20$ reward over $1{,}000$ pulls, and it never stops. The agent knows the
answer and keeps paying to re-confirm it.

That is the argument for decaying $\epsilon$ over time, and the argument for
UCB, which spends its exploration where uncertainty actually is.
\end{workedexample}

\end{v3topic}
\begin{v3topic}{Upper Confidence Bound (UCB)}
        The UCB1 algorithm and its regret bound are due to Auer,
        Cesa-Bianchi and
        Fischer\textsuperscript{\cite{auerFinitetimeAnalysisMultiarmed2002}}.

        Explore arms with high \emph{uncertainty} rather than randomly.
        \textbf{UCB1} selects:
        \begin{equation}
            A_t = \arg\max_a \left[Q_t(a) + c\,\sqrt{\frac{\ln t}{N_t(a)}}\right]
        \end{equation}
        where $N_t(a)$ is the number of times arm $a$ has been pulled.
        The confidence bonus decreases as an arm is explored more; $c$
        controls exploration width.
        UCB1 achieves $O(\ln T)$ regret --- asymptotically
        optimal.\textsuperscript{\cite[][p.~337]{hurbansGrokkingArtificialIntelligence2020}}

\begin{workedexample}[UCB deliberately picks the worse-looking arm]
At time $t = 12$, with $c = 2$:
arm~A has been pulled $10$ times with mean $0.6$; arm~B has been pulled twice
with mean $0.5$.

\[
\begin{aligned}
  \text{UCB}(A) &= 0.6 + 2\sqrt{\tfrac{\ln 12}{10}}
    = 0.6 + 2\sqrt{\tfrac{2.485}{10}}
    = 0.6 + 2(0.499) = 1.597,\\[3pt]
  \text{UCB}(B) &= 0.5 + 2\sqrt{\tfrac{\ln 12}{2}}
    = 0.5 + 2\sqrt{1.242}
    = 0.5 + 2(1.115) = 2.729 .
\end{aligned}
\]

UCB picks \textbf{B}, the arm with the \emph{lower} observed mean --- because
with only two samples its estimate could easily be wrong by more than the
$0.1$ gap. The bonus term $\sqrt{\ln t / N_t(a)}$ shrinks as $N_t(a)$ grows, so
once B has been pulled enough times its bonus collapses and the decision
reverts to the means.

This is the whole idea: \emph{optimism in the face of uncertainty}. Explore
what you do not know, not what a coin tells you to.
\end{workedexample}

\end{v3topic}
\begin{v3topic}{Thompson Sampling}
        A \textbf{Bayesian} approach: maintain a prior over arm reward
        parameters; sample from the posterior; pull the arm with the highest
        sampled value; update the posterior on the observed reward.

        For Bernoulli rewards, use a $\text{Beta}(\alpha_a, \beta_a)$ prior:
        \begin{itemize}
            \item Success: $\alpha_a \leftarrow \alpha_a + 1$
            \item Failure: $\beta_a \leftarrow \beta_a + 1$
        \end{itemize}
        Naturally balances exploration and exploitation via uncertainty.
        Empirically often outperforms $\epsilon$-greedy and
        UCB.\textsuperscript{\cite[][p.~338]{hurbansGrokkingArtificialIntelligence2020}}

\end{v3topic}


    % -------------------------------------------------------
    \section{Q-Learning and Deep Reinforcement Learning}


\begin{v3topic}{Monte Carlo Methods}
        \keyterm{Monte Carlo} methods estimate values by \emph{averaging
        complete returns}. They need no model of $P$ or $R$ and make no
        bootstrapping assumption --- but they can only learn from episodes
        that terminate\textsuperscript{\cite[][ch.~5]{suttonReinforcementLearningIntroduction2018}}.

        \textbf{First-visit MC prediction}: for each state $s$, average the
        returns $G_t$ observed after the first visit to $s$ in each episode:
        \begin{equation}
            V(s) \leftarrow \operatorname{average}\bigl(G_t \mid
            S_t = s,\ \text{first visit}\bigr).
            \label{eq:rl-mc-prediction}
        \end{equation}
        The estimate is unbiased and converges to $V^{\pi}(s)$ as the number
        of visits grows.

        \textbf{MC control} alternates policy evaluation with greedy policy
        improvement, using \emph{exploring starts} or an $\epsilon$-soft
        policy to guarantee that every state--action pair keeps being visited.

\begin{workedexample}[Monte Carlo against temporal difference on one episode]
An episode visits $A \to B \to$ terminal, receiving reward $0$ then $1$, with
$\gamma = 1$. Current estimates are $V(A) = 0.5$ and $V(B) = 0.5$; take
$\alpha = 0.1$.

\emph{Monte Carlo} waits for the episode to end and uses the actual return
from each state: $G_A = 0 + 1 = 1$ and $G_B = 1$. So
\[
  V(A) \leftarrow 0.5 + 0.1(1 - 0.5) = 0.55,
  \qquad
  V(B) \leftarrow 0.5 + 0.1(1 - 0.5) = 0.55 .
\]

\emph{TD(0)} updates $A$ immediately, using its \emph{current estimate} of
$B$ rather than waiting:
\[
  V(A) \leftarrow 0.5 + 0.1\bigl(0 + 1(0.5) - 0.5\bigr) = 0.5 \ \text{(no change)},
\]
then on the next transition
$V(B) \leftarrow 0.5 + 0.1(1 + 0 - 0.5) = 0.55$.

\emph{Read the difference.} MC moved $A$ on the strength of what actually
happened; TD did not move $A$ at all, because its estimate of $B$ already
agreed with the observed step. MC is unbiased but high variance --- the whole
noisy return enters the update. TD is biased --- it trusts its own estimate ---
but low variance, and it can learn from incomplete episodes and from
continuing tasks that never terminate. That bias--variance trade is the axis
along which every method in this chapter sits, and $n$-step methods and
TD($\lambda$) interpolate along it.
\end{workedexample}

\end{v3topic}
\begin{v3topic}{Temporal Difference Learning}
        \textbf{TD learning} is \emph{model-free}: learns directly from
        sampled transitions without a model of $P$ or $R$, and updates
        \emph{online} after each step.

        \textbf{TD(0)} value update:
        \begin{equation}
            V(s_t) \leftarrow V(s_t) + \alpha
            \underbrace{\bigl[r_{t+1} + \gamma V(s_{t+1}) - V(s_t)\bigr]}_{\text{TD error }\delta_t}
        \end{equation}
        TD(0) bootstraps from the current estimate; $\alpha$ is the learning
        rate.\textsuperscript{\cite[][p.~343]{hurbansGrokkingArtificialIntelligence2020}}

\end{v3topic}
\begin{v3topic}{SARSA vs.\ Q-Learning}
        Both learn $Q(s,a)$ via TD updates but differ in the \emph{target}:

        \textbf{SARSA} (on-policy):
        \begin{equation}
            Q(s,a) \leftarrow Q(s,a) + \alpha\bigl[r + \gamma Q(s',a') - Q(s,a)\bigr]
        \end{equation}
        Uses the \emph{actually taken} next action $a'$.

        \textbf{Q-Learning} (off-policy):
        \begin{equation}
            Q(s,a) \leftarrow Q(s,a) + \alpha\bigl[r + \gamma \max_{a'} Q(s',a') - Q(s,a)\bigr]
        \end{equation}
        Always bootstraps from the \emph{greedy} action; converges to
        $Q^*$ regardless of behaviour
        policy.\textsuperscript{\cite[][p.~346]{hurbansGrokkingArtificialIntelligence2020}}

\begin{workedexample}[One Q-learning update, term by term]
Current estimate $Q(s,a) = 2.0$. The agent acts, receives $r = 1$, and lands in
$s'$ where the best action is worth $\max_{a'} Q(s',a') = 5.0$. Take
$\gamma = 0.9$ and $\alpha = 0.1$.

\[
\begin{aligned}
  \text{TD target} &= r + \gamma\max_{a'}Q(s',a') = 1 + 0.9(5.0) = 5.5,\\
  \text{TD error } \delta &= 5.5 - 2.0 = 3.5,\\
  Q_{\text{new}}(s,a) &= 2.0 + 0.1(3.5) = \mathbf{2.35}.
\end{aligned}
\]

\emph{Read the terms.} The target $5.5$ is a \emph{better-informed guess} than
the old $2.0$, because it contains one real observed reward and only then falls
back on an estimate. The error $\delta = 3.5$ says the old estimate was
badly pessimistic. And $\alpha = 0.1$ means we move only a tenth of the way ---
the target is itself an estimate, so trusting it fully would make learning
unstable. Every value-based RL algorithm in this chapter is a variation on
these three lines.
\end{workedexample}

\begin{cautionbox}[On-policy and off-policy diverge exactly where it matters]
SARSA uses the action actually taken; Q-learning uses the best available
action. On a gridworld with a cliff along the shortest path, this produces
genuinely different behaviour: Q-learning learns the optimal path along the
cliff edge, because $\max_{a'}$ assumes it will never make an exploratory
mistake. SARSA learns a longer, safer path, because its target includes the
$\epsilon$-greedy chance of stepping off.

Neither is wrong. Q-learning answers ``what is best if I act perfectly?'';
SARSA answers ``what is best given that I sometimes explore?'' If your agent
will explore in the real environment --- a robot, a treatment policy, a trading
system --- SARSA is answering your question and Q-learning is not.
\end{cautionbox}

\end{v3topic}
\begin{v3topic}{Deep Q-Networks (DQN)}
        DQN is due to Mnih and
        colleagues\textsuperscript{\cite{mnihHumanlevelControlDeep2015}}; the
        tabular algorithm it generalises is due to Watkins and
        Dayan\textsuperscript{\cite{watkinsQlearning1992}}.

        For large or continuous state spaces, a Q-table is intractable.
        \textbf{DQN} approximates the Q-function with a deep neural network
        $Q(s,a;\theta) \approx Q^*(s,a)$, trained to minimise:
        \begin{equation}
            \mathcal{L}(\theta) = \mathbb{E}_{(s,a,r,s') \sim \mathcal{D}}
            \Bigl[\bigl(r + \gamma \max_{a'} Q(s',a';\theta^-) - Q(s,a;\theta)\bigr)^2\Bigr]
        \end{equation}
        Key innovations: (1)~\textbf{experience replay} buffer $\mathcal{D}$;
        (2)~\textbf{target network} $\theta^-$ updated periodically (not every
        step). Achieved human-level performance on 49 Atari
        games.\textsuperscript{\cite[][p.~350]{hurbansGrokkingArtificialIntelligence2020}}

\end{v3topic}
\begin{v3topic}{Experience Replay}
        Store each transition $(s, a, r, s')$ in a circular \textbf{replay
        buffer} $\mathcal{D}$; sample random \emph{mini-batches} to train.

        Benefits:
        \begin{enumerate}
            \item \textbf{Decorrelates} temporally adjacent samples, reducing
                  variance
            \item \textbf{Sample efficiency}: each transition reused many times
            \item \textbf{Stabilises} gradient updates by breaking
                  non-stationarity
        \end{enumerate}
        \textsuperscript{\cite[][p.~351]{hurbansGrokkingArtificialIntelligence2020}}

\end{v3topic}
\begin{v3topic}{Double Q-Learning}
        Standard Q-learning suffers from \textbf{overestimation bias}: the
        same network selects \emph{and} evaluates the maximum action, biasing
        targets upward.

        \textbf{Double DQN} decouples selection and evaluation:
        \begin{equation}
            y = r + \gamma\, Q\!\left(s',\, \arg\max_{a'} Q(s',a';\theta);\;\theta^-\right)
        \end{equation}
        Online network $\theta$ selects; target network $\theta^-$ evaluates.
        Reduces overestimation, yields more stable training and better
        policies.\textsuperscript{\cite[][p.~353]{hurbansGrokkingArtificialIntelligence2020}}

\end{v3topic}
\begin{v3topic}{Sparse Rewards and Reward Shaping}
        Many real tasks only give reward at task completion --- the agent faces
        a hard \textbf{credit assignment} problem.

        \textbf{Mitigation strategies}:
        \begin{itemize}
            \item \textbf{Reward shaping}: add potential-based intermediary
                  rewards $F(s,s') = \gamma \Phi(s') - \Phi(s)$ without
                  changing optimal policy
            \item \textbf{Intrinsic motivation}: curiosity bonus for novel states
            \item \textbf{Curriculum learning}: start with easier sub-tasks
            \item \textbf{Hindsight Experience Replay}: relabel failed
                  trajectories toward a different goal
        \end{itemize}
        \textsuperscript{\cite[][p.~357]{hurbansGrokkingArtificialIntelligence2020}}

\end{v3topic}
\begin{v3topic}{Hierarchical Reinforcement Learning}
        Decomposes a long-horizon task into \textbf{sub-goals} and
        \textbf{sub-policies} at multiple levels of temporal abstraction.

        \textbf{Options framework}: an option is a triple
        $(I, \pi_\omega, \beta)$:
        \begin{itemize}
            \item $I \subseteq \mathcal{S}$: initiation set
            \item $\pi_\omega(a|s)$: intra-option policy
            \item $\beta(s)$: termination probability
        \end{itemize}
        A high-level policy selects options; each option executes its
        sub-policy until termination. Enables reuse of learned
        skills.\textsuperscript{\cite[][p.~358]{hurbansGrokkingArtificialIntelligence2020}}

\end{v3topic}
\begin{v3topic}{Value-Based vs.\ Policy-Based Methods}
        \begin{tblr}{width=\linewidth,
            colspec={X[l]X[l]},
            hlines, vlines,
            row{1}={font=\bfseries\small, bg=LimeGreen!20},
            row{2-Z}={font=\small},
        }
            Value-Based & Policy-Based \\
            Learn $Q$ or $V$; derive $\pi$ & Learn $\pi$ directly \\
            Q-learning, DQN & REINFORCE, PPO \\
            Natural for discrete actions & Handles continuous actions \\
            Deterministic output & Stochastic policy \\
            Can be unstable & High variance gradients \\
        \end{tblr}
\captionof{table}{Value-based against policy-based methods. Value methods need a tractable maximisation over actions, which is why continuous action spaces usually call for policy methods.}
\label{tab:rl-value-policy}
        \textsuperscript{\cite[][p.~360]{hurbansGrokkingArtificialIntelligence2020}}

\end{v3topic}
\begin{v3topic}{Actor-Critic Methods}
        Combines value-based and policy-based approaches:
        \begin{itemize}
            \item \textbf{Actor} $\pi(a|s;\theta)$: parametric policy, updated
                  via gradient ascent
            \item \textbf{Critic} $V(s; w)$: estimates state values, reduces
                  gradient variance
        \end{itemize}
        The \textbf{advantage function} $A(s,a) = Q(s,a) - V(s)$ measures
        how much better action $a$ is than average, reducing variance.

        \textbf{A2C}: synchronous workers; \textbf{A3C}: asynchronous
        parallel workers eliminating need for a replay
        buffer.\textsuperscript{\cite[][p.~362]{hurbansGrokkingArtificialIntelligence2020}}

\end{v3topic}
    % -------------------------------------------------------
    \section{Reinforcement Learning Approaches}

    \subsection{Model-Free Learning}

\begin{v3topic}{Model-Free RL}
    In \textbf{model-free} reinforcement learning the agent learns entirely through
    interaction with the environment --- no internal model of state transitions or
    rewards is built. The agent discovers which actions are valuable solely from
    trial-and-error experience.\textsuperscript{\cite[][p.~348]{hurbansGrokkingArtificialIntelligence2020}}

    Key characteristics:
    \begin{itemize}
        \item No explicit environment model required
        \item Learns from $(s, a, r, s’)$ tuples directly
        \item More data-hungry than model-based, but robust to model errors
        \item Examples: Q-learning, SARSA, REINFORCE
    \end{itemize}

\end{v3topic}
\begin{v3topic}{The Q-Learning Update}
    \textbf{Q-learning} learns the action-value function $Q(s,a)$ --- the expected
    cumulative reward from taking action $a$ in state $s$ and following the
    optimal policy thereafter.\textsuperscript{\cite[][pp.~335--340]{hurbansGrokkingArtificialIntelligence2020}}

    The \textbf{Bellman update rule}:
    \begin{equation}
        Q(s,a) \leftarrow (1{-}\alpha)\,Q(s,a)
                 + \alpha\!\left[r + \gamma \max_{a’} Q(s’,a’)\right]
    \end{equation}
    $\alpha$: learning rate; $\gamma$: discount factor; $r$: immediate reward;
    $s’$: next state.

    The Q-table (rows = states, columns = actions) is initialised to zero and
    updated incrementally. At inference time: $a^* = \arg\max_a Q(s,a)$.

\end{v3topic}
\begin{v3topic}{SARSA: On-Policy TD Control}
    \textbf{SARSA} is \emph{on-policy}: it updates $Q$ using the action $a’$
    \emph{actually taken} by the current policy, not the greedy
    maximum.\textsuperscript{\cite[][Ch.~6]{suttonReinforcementLearningIntroduction2018}}

    \begin{equation}
        Q(s,a) \leftarrow Q(s,a) + \alpha\bigl[r + \gamma\,Q(s’,a’) - Q(s,a)\bigr]
    \end{equation}

    \begin{itemize}
        \item On-policy: more conservative, avoids dangerous exploratory actions
        \item Q-learning: faster convergence, may take risky shortcuts in training
        \item Both converge to the optimal $Q^*$ given sufficient exploration
    \end{itemize}

\end{v3topic}
\begin{v3topic}{REINFORCE (Policy Gradient)}
    \textbf{REINFORCE} directly optimises the policy $\pi_\theta(a|s)$ by
    gradient ascent on expected
    return.\textsuperscript{\cite[][Ch.~13]{suttonReinforcementLearningIntroduction2018}}

    \begin{equation}
        \nabla_\theta J(\theta) = \mathbb{E}_\pi\!\left[
            \nabla_\theta \ln \pi_\theta(a_t|s_t) \cdot G_t\right]
    \end{equation}

    where $G_t = \sum_{k=0}^{\infty} \gamma^k r_{t+k+1}$ is the return from step $t$.

    \begin{itemize}
        \item Monte Carlo: updates after complete episodes
        \item High variance; baseline subtraction (actor-critic) reduces variance
        \item Natural for continuous action spaces
    \end{itemize}

\end{v3topic}

    \subsection{Model-Based Learning}

\begin{v3topic}{Model-Based RL}
    In \textbf{model-based} RL the agent learns or is given an explicit model
    $\hat{M}(s,a) \to (s’, r)$ of the environment. It can \emph{plan} by
    simulating experience inside the model without executing every
    action.\textsuperscript{\cite[][p.~348]{hurbansGrokkingArtificialIntelligence2020}}

    \begin{itemize}
        \item \textbf{Sample efficient}: plan many steps from few real interactions
        \item Can reason about counterfactual actions before committing
        \item Model errors compound over long horizons --- accuracy matters
        \item Trade-off: model capacity vs.\ model bias
    \end{itemize}

\end{v3topic}
\begin{v3topic}{Dyna-Q: Planning + Learning}
    \textbf{Dyna-Q} (Sutton, 1991) combines real experience with simulated
    experience from a learned
    model:\textsuperscript{\cite[][Ch.~8]{suttonReinforcementLearningIntroduction2018}}

    \begin{enumerate}
        \item Take action $a$; observe $r$, $s’$ from real environment
        \item Update Q-table with real transition
        \item Update model: $\hat{M}(s,a) \leftarrow (s’,r)$
        \item Repeat $n$ times: sample random past $(s,a)$, simulate
              $(s’,r)$ from model, update Q
    \end{enumerate}

    The $n$ planning steps give Dyna-Q major sample-efficiency gains.
    Planning in imagination is cheap; real experience is expensive.

\end{v3topic}
\begin{v3topic}{World Models \& MPC}
    \textbf{World models} (Ha \& Schmidhuber, 2018) learn a compact
    latent representation of the environment; the agent plans entirely within
    the latent model.\textsuperscript{\cite[][Ch.~8]{suttonReinforcementLearningIntroduction2018}}

    \textbf{Model Predictive Control (MPC)}: at each step, optimise actions
    over a finite horizon $H$ using the model, execute only the first action,
    then re-plan. Naturally handles non-stationary environments.

    \begin{itemize}
        \item Used in robotics, process control, AlphaZero (MCTS + learned model)
        \item Horizon $H$ trades off computation vs.\ look-ahead quality
    \end{itemize}

\end{v3topic}

    \subsection{Hierarchical Reinforcement Learning}

\begin{v3topic}{Temporal Abstraction and Options}
        Flat RL chooses a primitive action every step, which makes long tasks
        with sparse rewards effectively unlearnable --- the credit-assignment
        chain is simply too long (see the value-decay calculation earlier in
        this chapter). \keyterm{Hierarchical reinforcement learning} attacks
        this by letting the agent choose over \emph{temporally extended}
        behaviours rather than single actions.

        The \keyterm{options framework}\textsuperscript{\cite{suttonBetweenMDPsSemiMDPs1999}}
        formalises an option as a triple
        \begin{equation}
            \omega = \bigl(\mathcal{I}_\omega,\ \pi_\omega,\
            \beta_\omega\bigr),
            \label{eq:rl-option}
        \end{equation}
        where $\mathcal{I}_\omega \subseteq \mathcal{S}$ is the set of states
        in which the option may be started, $\pi_\omega$ is the policy it
        follows while running, and $\beta_\omega(s) \in [0,1]$ is the
        probability of terminating in state $s$. A primitive action is the
        special case of an option that always terminates after one step.

        Because options take variable time, the decision problem over options
        is a \keyterm{semi-Markov decision process} rather than an MDP, and
        the Bellman equations carry a discount raised to the option's random
        duration.

        \textbf{MAXQ}\textsuperscript{\cite{dietterichHierarchicalReinforcementLearning2000}}
        takes a complementary route: it decomposes the value function itself
        along a given task hierarchy, so that sub-task value functions can be
        learned once and reused wherever the sub-task appears.

\begin{plainlanguage}
Flat reinforcement learning is like navigating a city by deciding which way to
face at every single footstep. Hierarchical RL lets the agent decide ``go to
the station'', hand control to a sub-policy that knows how to walk there, and
only make the next high-level decision on arrival. The rewards are the same;
the number of decisions between action and consequence collapses.
\end{plainlanguage}

\begin{workedexample}[Why abstraction rescues a sparse reward]
A goal lies $100$ primitive steps away and pays $+1$ on arrival, with
$\gamma = 0.9$. At the start state the value is
\[
  V \approx \gamma^{99} = 0.9^{99} = 3.0\times10^{-5},
\]
numerically indistinguishable from zero, so the gradient of the start-state
value with respect to anything is negligible.

Now suppose four options each cover about $25$ primitive steps. At the option
level the goal is $4$ decisions away. If the discount is applied per
\emph{decision} rather than per step, the start-state value becomes
\[
  \gamma^{3} = 0.9^{3} = 0.729 ,
\]
a signal four orders of magnitude larger. Nothing about the environment or the
reward changed --- only the granularity at which credit is assigned.

\emph{The catch.} The options must come from somewhere. Hand-designing them
requires the domain knowledge you were hoping to learn, and discovering them
automatically remains an open research problem. That gap is why hierarchical
RL is conceptually compelling and comparatively rare in production.
\end{workedexample}

\end{v3topic}

    \subsection{Exploration vs Exploitation}

\begin{v3topic}{The Exploration--Exploitation Trade-off}
    The agent must balance:\textsuperscript{\cite[][pp.~338--339]{hurbansGrokkingArtificialIntelligence2020}}
    \begin{itemize}
        \item \textbf{Exploitation}: select $\arg\max_a Q(s,a)$ to maximise
              known reward
        \item \textbf{Exploration}: try other actions to discover potentially
              better rewards
    \end{itemize}

    Purely greedy agents get stuck in local optima; purely random agents
    never converge. Every RL algorithm must address this tension.

\end{v3topic}
\begin{v3topic}{$\varepsilon$-Greedy Strategy}
    The simplest trade-off
    strategy:\textsuperscript{\cite[][p.~339]{hurbansGrokkingArtificialIntelligence2020}}

    \begin{equation}
        a_t = \begin{cases}
            \text{random action} & \text{with prob.\ } \varepsilon \\
            \arg\max_a Q(s_t,a) & \text{with prob.\ } 1{-}\varepsilon
        \end{cases}
    \end{equation}

    \textbf{Decaying $\varepsilon$}: start at $\varepsilon \approx 1.0$ (mostly
    explore) and anneal to $\varepsilon \approx 0.01$ (mostly exploit) over
    training:
    \[ \varepsilon_t = \varepsilon_{\min}
       + (\varepsilon_{\max} - \varepsilon_{\min})\,e^{-t/\tau} \]

\end{v3topic}
\begin{v3topic}{Upper Confidence Bound (UCB)}
    UCB selects actions that are either \emph{estimated good} or
    \emph{rarely tried}:\textsuperscript{\cite[][Ch.~2]{suttonReinforcementLearningIntroduction2018}}

    \begin{equation}
        a_t = \arg\max_a \left[ Q(s,a) + c\sqrt{\frac{\ln t}{N_t(s,a)}} \right]
    \end{equation}

    $N_t(s,a)$ = times action $a$ was taken from state $s$; $c$ controls
    exploration strength. The confidence bonus shrinks as an action is tried
    more, automatically directing attention to underexplored options.
    \textbf{Optimism in the face of uncertainty.}

\end{v3topic}
\begin{v3topic}{Curiosity-Driven Exploration}
    When extrinsic rewards are sparse, \textbf{intrinsic motivation} sustains
    exploration. A \emph{forward model} $f(s_t, a_t) \to \hat{s}_{t+1}$ is
    learned; its prediction error provides an intrinsic
    reward:\textsuperscript{\cite[][Ch.~2]{suttonReinforcementLearningIntroduction2018}}

    \begin{equation}
        r^{\mathrm{int}}_t = \bigl\lVert\hat{s}_{t+1} - s_{t+1}\bigr\rVert^2
    \end{equation}

    High error signals novel, informative states the model has not yet learned.
    The agent is intrinsically \emph{curious} about regions where its world
    model is poor. Successfully applied in Montezuma’s Revenge and robotic
    manipulation.

\end{v3topic}


    % -------------------------------------------------------
    \section{Inference and Causality}

    \subsection{Statistical Inference}

\begin{v3topic}{Bayesian Inference}
    Bayesian inference treats model parameters $\theta$ as random variables
    and updates beliefs upon observing data
    $\mathcal{D}$:\textsuperscript{\cite[][pp.~14--19]{barberBayesianReasoningMachine2012}}

    \begin{equation}
        p(\theta|\mathcal{D}) = \frac{p(\mathcal{D}|\theta)\,p(\theta)}{p(\mathcal{D})}
    \end{equation}

    \begin{itemize}
        \item \textbf{Prior} $p(\theta)$: belief before seeing data
        \item \textbf{Likelihood} $p(\mathcal{D}|\theta)$: how probable is the data?
        \item \textbf{Posterior} $p(\theta|\mathcal{D})$: updated belief
        \item \textbf{Evidence} $p(\mathcal{D}) = \int p(\mathcal{D}|\theta)\,p(\theta)\,d\theta$
    \end{itemize}

    Posterior $\propto$ Likelihood $\times$ Prior.
    The MAP estimate is $\hat{\theta} = \arg\max_\theta p(\theta|\mathcal{D})$.

\end{v3topic}
\begin{v3topic}{Bayesian vs Frequentist}
    \begin{itemize}
        \item \textbf{Frequentist}: $\theta$ is a fixed unknown constant;
              probability refers to long-run frequencies; provides point
              estimates and confidence intervals
        \item \textbf{Bayesian}: $\theta$ is a random variable with a
              distribution; probability reflects degree of belief; provides
              full posterior distribution
    \end{itemize}

    Bayesian inference is naturally sequential: today’s posterior becomes
    tomorrow’s prior as new data
    arrives.\textsuperscript{\cite[][Ch.~1]{downeyThinkBayes2016}}
    Ideal for online learning and decision-making under uncertainty.

\end{v3topic}
\begin{v3topic}{Bayesian Networks}
    A \textbf{Bayesian Network} (belief network) is a DAG where each node is
    a random variable $x_i$ with a conditional probability table (CPT)
    $p(x_i | \mathrm{pa}(x_i))$. The joint distribution
    factorises:\textsuperscript{\cite[][pp.~33--40]{barberBayesianReasoningMachine2012}}

    \begin{equation}
        p(x_1, \ldots, x_D) = \prod_{i=1}^{D} p\bigl(x_i \mid \mathrm{pa}(x_i)\bigr)
    \end{equation}

    \textbf{Wet Grass example}: $R$ (rain), $S$ (sprinkler), $T$ (Tracey’s
    grass), $J$ (Jack’s grass). With $p(T|J,R,S)=p(T|R,S)$ and $p(J|R,S)=p(J|R)$:
    \[p(T,J,R,S) = p(T|R,S)\,p(J|R)\,p(R)\,p(S)\]
    This compresses 15 parameters to 8. The graph encodes sets of conditional
    independence relations through d-separation and the local Markov property;
    an absent edge should not be read as one unconditional independence claim.
\visualseparator
    \textbf{Explaining away}: in $A \to C \leftarrow B$, causes $A$ and $B$
    are independent a priori. Conditioning on the effect $C$ makes them
    dependent: observing one cause ``explains away’’ the other.
    This is the collider effect.

\end{v3topic}
\begin{v3topic}{Probabilistic Modelling}
    A probabilistic model defines a joint distribution over observed variables
    $\mathbf{x}$ and latent variables
    $\mathbf{z}$:\textsuperscript{\cite[][Ch.~1]{barberBayesianReasoningMachine2012}}

    \begin{equation}
        p(\mathbf{x}, \mathbf{z} | \theta) =
            p(\mathbf{x} | \mathbf{z}, \theta)\,p(\mathbf{z}|\theta)
    \end{equation}

    \textbf{Inference}: compute $p(\mathbf{z}|\mathbf{x},\theta)$ given observations.\\
    \textbf{Learning}: find $\theta$ maximising the marginal likelihood
    $p(\mathbf{x}|\theta) = \int p(\mathbf{x}|\mathbf{z},\theta)\,p(\mathbf{z}|\theta)\,d\mathbf{z}$.

    The integral is often intractable, motivating approximate methods:
    MCMC and variational inference.

\end{v3topic}
\begin{v3topic}{Conditional Independence}
    $X$ and $Y$ are \textbf{conditionally independent} given $Z$, written
    $X \perp\!\!\!\perp Y \mid Z$,
    if:\textsuperscript{\cite[][p.~8]{barberBayesianReasoningMachine2012}}

    \begin{equation}
        p(X, Y | Z) = p(X|Z)\,p(Y|Z)
    \end{equation}

    Equivalently, $p(X|Y,Z) = p(X|Z)$: once $Z$ is known, $Y$ contains no
    additional information about $X$.

    Every absent edge in a Bayesian Network encodes exactly one conditional
    independence assumption, compressing the representation of the joint
    distribution.

\end{v3topic}

    \subsection{Introduction to Causality}

\begin{v3topic}{Correlation vs Causation}
    \textbf{Correlation} is a symmetric statistical association;
    $\mathrm{Cov}(X,Y) \neq 0$ says nothing about direction or mechanism.

    \textbf{Causation} implies a directional mechanism: intervening on $X$
    forces a change in $Y$.

    Classic example: ice cream sales and drowning rates are correlated ---
    both are driven by a common cause (summer heat). Controlling for
    temperature removes the association. Formally: $p(Y|X) \neq p(Y|\,\mathrm{do}(X))$
    whenever confounders
    exist.\textsuperscript{\cite[][Ch.~1]{pearlBookWhyNew2018}}

\end{v3topic}
\begin{v3topic}{Granger Causality}
    \textbf{Granger causality} (Granger, 1969): $X$ Granger-causes $Y$ if past
    values of $X$ provide statistically significant predictive information
    about $Y$, beyond what past $Y$ alone
    provides:\textsuperscript{\cite[][Ch.~2]{nessCausalAI2025}}

    \[ Y_t = f(Y_{t-1}, \ldots, X_{t-1}, \ldots) + \varepsilon_t \]

    Test: does adding $X_{t-k}$ significantly reduce prediction error for $Y_t$?

    \textbf{Limitation}: captures temporal precedence, not structural
    causation. Two variables driven by a hidden common cause can each appear
    to Granger-cause the other.

\end{v3topic}
\begin{v3topic}{Directed Acyclic Graphs (DAGs)}
    A \textbf{DAG} encodes causal structure: directed edges $X \to Y$ mean
    ``$X$ directly causes $Y$’’. Acyclicity prevents circular
    causation.\textsuperscript{\cite[][pp.~25--27]{barberBayesianReasoningMachine2012}}

    DAG terminology:
    \begin{itemize}
        \item \textbf{Parent}: direct cause, $\mathrm{pa}(Y)$
        \item \textbf{Ancestor}: any cause along a directed path
        \item \textbf{Descendant}: any variable causally downstream
        \item \textbf{Markov blanket}: parents + children + parents of children;
              screens $X$ from all other variables
    \end{itemize}

    In causal DAGs, ancestral order corresponds to temporal order.

\end{v3topic}
\begin{v3topic}{Fork, Chain, and Collider}
    Three elementary structures determine how information
    flows:\textsuperscript{\cite[][pp.~43--45]{barberBayesianReasoningMachine2012}}

    \begin{itemize}
        \item \textbf{Fork} (common cause): $A \leftarrow C \rightarrow B$.
              $C$ is a confounder; $A \perp\!\!\!\perp B \mid C$ but not marginally.
        \item \textbf{Chain} (mediation): $A \rightarrow C \rightarrow B$.
              $A \perp\!\!\!\perp B \mid C$ but not marginally.
        \item \textbf{Collider}: $A \rightarrow C \leftarrow B$.
              $A \perp\!\!\!\perp B$ marginally, but $A \not\!\perp\!\!\!\perp B \mid C$.
              Conditioning opens the path.
    \end{itemize}

    Conditioning on a collider (or its descendant) creates a \emph{spurious}
    association between its parents.

\end{v3topic}
\begin{v3topic}{D-Separation}
    \textbf{D-separation} (Pearl, 1988) is the formal criterion for reading
    conditional independence from a
    DAG.\textsuperscript{\cite[][pp.~45--46]{barberBayesianReasoningMachine2012}}

    A path $U$ between $X$ and $Y$ is \textbf{blocked} by set $\mathcal{Z}$ if:
    \begin{enumerate}
        \item There is a \emph{non-collider} $w$ on $U$ with $w \in \mathcal{Z}$, OR
        \item There is a \emph{collider} $w$ on $U$ such that neither $w$ nor
              any descendant of $w$ is in $\mathcal{Z}$.
    \end{enumerate}

    $X$ and $Y$ are \textbf{d-separated} by $\mathcal{Z}$ if every path between
    them is blocked. D-separation implies $X \perp\!\!\!\perp Y \mid \mathcal{Z}$
    in any distribution consistent with the DAG structure.
\visualseparator
    \textbf{Intuition}: d-separation tells us which conditioning sets render
    variables independent. It is the key tool for identifying valid adjustment
    sets without running experiments.

\end{v3topic}

    \subsection{Interventions}

\begin{v3topic}{Seeing vs Doing: The Do-Operator}
    Pearl’s \textbf{do-operator} distinguishes observation from
    intervention:\textsuperscript{\cite[][pp.~50--51]{barberBayesianReasoningMachine2012}}

    \begin{itemize}
        \item \textbf{Seeing}: $p(Y \mid X{=}x)$ --- conditioning on
              \emph{observing} $X{=}x$; confounders still act
        \item \textbf{Doing}: $p(Y \mid \mathrm{do}(X{=}x))$ --- \emph{forcing}
              $X{=}x$ by intervention; removes all incoming arrows to $X$
    \end{itemize}

    Intervening creates modified graph $\mathcal{G}_{\overline{X}}$ in which
    all edges \emph{into} $X$ are cut. Post-intervention: only $p(X_j | \mathrm{pa}(X_j))$
    for non-intervened variables remains; the intervened variable’s CPT is replaced
    by the unit distribution at the set value.

\end{v3topic}
\begin{v3topic}{Confounders}
    A \textbf{confounder} is a variable $C$ that causally affects both
    treatment $X$ and outcome
    $Y$:\textsuperscript{\cite[][Ch.~7]{hernanCausalInferenceWhat2020}}

    \[ C \rightarrow X,\quad C \rightarrow Y \]

    Confounders create a \emph{backdoor path} from $X$ to $Y$ that does not
    represent the causal effect of $X$. Failing to adjust for confounders
    yields biased estimates.

    \textbf{Examples}: socioeconomic status confounds education and health
    outcomes; age confounds treatment choice and recovery in clinical data.
    In a DAG, confounders are identified by backdoor paths.

\end{v3topic}
\begin{v3topic}{Counterfactuals}
    A \textbf{counterfactual} asks: ``What would $Y$ have been if $X$ had been
    $x$, given that we actually observed $X{=}x’$?’’\textsuperscript{\cite[][Ch.~7]{pearlCausalityModelsReasoning2009}}

    Notation: $Y_{X=x}(u)$ = value of $Y$ for unit $u$ under
    $\mathrm{do}(X{=}x)$.

    Requires a \textbf{Structural Causal Model} (SCM):
    $X_i = f_i(\mathrm{pa}(X_i), U_i)$.
    Three steps:
    \begin{enumerate}
        \item \textbf{Abduction}: infer exogenous noise $U$ from evidence
        \item \textbf{Action}: modify SCM to set $X{=}x$
        \item \textbf{Prediction}: compute $Y$ in the modified model
    \end{enumerate}

\end{v3topic}
\begin{v3topic}{Causal Inference vs RCTs}
    A \textbf{Randomized Controlled Trial (RCT)} randomly assigns treatment,
    breaking the confounder--treatment
    link:\textsuperscript{\cite[][Ch.~1]{hernanCausalInferenceWhat2020}}

    \[ \text{Random assignment} \;\Leftrightarrow\; \mathrm{do}(X)
       \;\Rightarrow\; p(Y|\mathrm{do}(X)) = p(Y|X) \]

    RCTs are the gold standard but often costly, unethical, or impossible.

    \textbf{Observational causal inference} estimates $p(Y|\mathrm{do}(X))$
    from non-experimental data via:
    \begin{itemize}
        \item Backdoor/front-door adjustment (Pearl’s criteria)
        \item Instrumental variables, regression discontinuity,
              difference-in-differences
    \end{itemize}

\end{v3topic}

    \subsection{Do-Calculus}

\begin{v3topic}{Backdoor Criterion}
    A set $\mathcal{Z}$ satisfies the \textbf{backdoor criterion} for
    $(X,Y)$ if:\textsuperscript{\cite[][pp.~79--81]{pearlCausalityModelsReasoning2009}}
    \begin{enumerate}
        \item No node in $\mathcal{Z}$ is a descendant of $X$
        \item $\mathcal{Z}$ blocks every backdoor path from $X$ to $Y$
    \end{enumerate}

    If satisfied, the \textbf{adjustment formula} gives:
    \begin{equation}
        p(Y \mid \mathrm{do}(X)) = \sum_{z} p(Y \mid X, \mathcal{Z}{=}z)\,p(\mathcal{Z}{=}z)
    \end{equation}

    This allows computing causal effects from observational data by conditioning
    on the confounder set $\mathcal{Z}$.

\end{v3topic}
\begin{v3topic}{Front-Door Criterion}
    Used when all confounders are \emph{unobserved} but a mediator $M$ is
    observed:\textsuperscript{\cite[][pp.~82--84]{pearlCausalityModelsReasoning2009}}

    Conditions: (1) $M$ intercepts all directed paths $X \to Y$; (2) no
    unblocked backdoor path from $X$ to $M$; (3) all backdoor paths from $M$
    to $Y$ are blocked by $X$.

    \begin{equation}
        p(Y|\mathrm{do}(X)) = \sum_m p(M{=}m|X)
            \sum_{x’} p(Y|M{=}m, X{=}x’)\,p(X{=}x’)
    \end{equation}

\end{v3topic}
\begin{v3topic}{Three Rules of Do-Calculus}
    Pearl’s \textbf{do-calculus} provides three rules that together are
    \emph{complete}: any identifiable causal query can be derived from
    observational data + DAG using these
    transformations.\textsuperscript{\cite[][pp.~85--87]{pearlCausalityModelsReasoning2009}}

    Let $\mathcal{G}_{\overline{X}}$ denote graph with arrows into $X$ removed;
    $\mathcal{G}_{\underline{X}}$ graph with arrows out of $X$ removed.

    \begin{itemize}
        \item \textbf{Rule 1 (Observation)} if $(Y \perp\!\!\!\perp Z \mid X,W)_{\mathcal{G}_{\overline{X}}}$:
            $p(y|\mathrm{do}(x),z,w) = p(y|\mathrm{do}(x),w)$
        \item \textbf{Rule 2 (Action/observation exchange)} if
            $(Y \perp\!\!\!\perp Z \mid X,W)_{\mathcal{G}_{\overline{X}\underline{Z}}}$:
            $p(y|\mathrm{do}(x),\mathrm{do}(z),w) = p(y|\mathrm{do}(x),z,w)$
        \item \textbf{Rule 3 (Action deletion)} if
            $(Y \perp\!\!\!\perp Z \mid X,W)_{\mathcal{G}_{\overline{X},\overline{Z(W)}}}$:
            $p(y|\mathrm{do}(x),\mathrm{do}(z),w) = p(y|\mathrm{do}(x),w)$
    \end{itemize}

\end{v3topic}

    \subsection{Fallacies}

\begin{v3topic}{Simpson’s Paradox}
    A statistical trend in \emph{combined} data reverses in every
    subgroup.\textsuperscript{\cite[][pp.~49--50]{barberBayesianReasoningMachine2012}}

    \textbf{Medical trial example}: A drug appears beneficial overall
    (recovery: 50\% drug vs.\ 40\% no-drug) but is \emph{harmful} in each
    subgroup (males: 60\% vs.\ 70\%; females: 20\% vs.\ 30\%).

\begin{workedexample}[Simpson's paradox in real clinical data]
Two treatments for kidney stones, from a study of $700$
patients\textsuperscript{\cite{charigComparisonTreatmentRenal1986}}:

\smallskip
\begin{tblr}{
  width=\linewidth,
  colspec={Q[l,0.22\linewidth] Q[c] Q[c] Q[c]},
  row{1}={font=\bfseries,bg=AlmanachMist},
  hlines,vlines}
& Small stones & Large stones & \textbf{Overall} \\
Treatment A & $81/87 = \mathbf{93\,\%}$ & $192/263 = \mathbf{73\,\%}$ &
  $273/350 = 78\,\%$ \\
Treatment B & $234/270 = 87\,\%$ & $55/80 = 69\,\%$ &
  $289/350 = \mathbf{83\,\%}$ \\
\end{tblr}
\smallskip

Read the table twice. Treatment~A wins on small stones ($93$ vs $87$). Treatment
A wins on large stones ($73$ vs $69$). Treatment~A \emph{loses} overall
($78$ vs $83$).

\emph{This is not an arithmetic error.} Check it: $81+192 = 273$ out of
$87+263 = 350$, and $234+55 = 289$ out of $270+80 = 350$. Both columns are
correct.

\emph{The resolution is causal, not statistical.} Look at the group sizes.
Treatment~A was given mostly to \emph{large}-stone patients ($263$ of $350$),
which are the harder cases; treatment~B mostly to small-stone patients ($270$
of $350$). Stone size affects both which treatment was chosen and whether the
patient recovers --- it is a confounder. The overall column silently compares a
hard caseload against an easy one.

\emph{Which number should a doctor use?} The subgroup ones. The correct
quantity is $p(\text{recovery}\mid \mathrm{do}(\text{treatment}))$, obtained by
adjusting for stone size, not $p(\text{recovery}\mid\text{treatment})$ read off
the aggregate. No amount of extra data fixes the aggregate figure, because the
problem is not sampling error --- it is that the question being answered is the
wrong one.
\end{workedexample}

\begin{cautionbox}[You cannot detect this from the aggregate alone]
Nothing in the overall column signals that it is misleading. The only defence
is to know, in advance and from domain knowledge, which variables might
influence both the treatment and the outcome --- and then to check whether the
effect survives conditioning on them. This is why a causal diagram is drawn
before the analysis, not after. It is also why ``we have millions of rows'' is
no protection at all: more data makes a confounded estimate more precise, not
less wrong.
\end{cautionbox}

    \textbf{Resolution}: The paradox conflates $p(R|D)$ (observational)
    with $p(R|\mathrm{do}(D))$ (interventional). Gender is a confounder
    affecting both drug choice and recovery. Applying the backdoor
    adjustment reveals the drug has no net benefit.

\end{v3topic}
\begin{v3topic}{Collider Bias \& Berkson’s Paradox}
    \textbf{Collider bias}: conditioning on a common effect of two variables
    creates a spurious association between
    them.\textsuperscript{\cite[][p.~43]{barberBayesianReasoningMachine2012}}

    Structure: $A \to C \leftarrow B$. Marginally $A \perp\!\!\!\perp B$, but
    $A \not\!\perp\!\!\!\perp B \mid C$.

    \textbf{Berkson’s Paradox} (1946): In hospital studies, diseases $A$ and
    $B$ appear negatively correlated because hospitalisation is a collider.
    Patients with $A$ seem less likely to also have $B$ --- but only because
    both independently cause hospitalisation. In the general population, $A$
    and $B$ are independent.

\end{v3topic}
\begin{v3topic}{Mediation Fallacy}
    The \textbf{mediation fallacy} occurs when a researcher adjusts for a
    mediator $M$ on the causal path $X \to M \to Y$ while intending to
    estimate the total effect of $X$ on
    $Y$.\textsuperscript{\cite[][Ch.~9]{pearlCausalityModelsReasoning2009}}

    Adjusting for $M$:
    \begin{itemize}
        \item \emph{Blocks} the indirect pathway $X \to M \to Y$
        \item The estimate captures only the \emph{direct} effect $X \to Y$,
              not the total effect
        \item If the goal is the total effect, $M$ must \emph{not} be
              conditioned on
    \end{itemize}

    Correct mediation analysis separates direct and indirect effects using
    counterfactual definitions.

\end{v3topic}
\begin{v3topic}{Missing Values: Causal View}
    Standard imputation (mean/median/model) assumes data is Missing At
    Random (MAR). The causal view requires modelling the
    \textbf{missingness mechanism}:\textsuperscript{\cite[][Ch.~8]{nessCausalAI2025}}

    \begin{itemize}
        \item \textbf{MCAR} (Missing Completely At Random): missingness
              independent of all variables --- safe to impute
        \item \textbf{MAR} (Missing At Random): missingness depends on
              observed variables --- adjust for them
        \item \textbf{MNAR} (Missing Not At Random): missingness depends
              on the missing value itself --- selection bias; model the
              mechanism explicitly
    \end{itemize}

    If the variable with missing values is a collider descendant, naive
    imputation can introduce collider bias. The DAG determines the
    correct imputation strategy.

\end{v3topic}

    % -------------------------------------------------------
\section{Deep Dive: Offline, Safe, and Preference-Based Learning}
\begin{v3topic}{Offline RL and Coverage}
Offline RL learns from a fixed behaviour dataset.
Actions outside that dataset's support can receive optimistic values because
their outcomes are weakly identified.
Report behaviour-policy coverage and compare with behaviour cloning and a
conservative non-learning baseline.

\end{v3topic}
\begin{v3topic}{Safety Is a Constraint}
Reward penalties do not guarantee a hard safety limit.
Constrained policies, shields, action filters, simulators, and independent
supervisors address different layers.
Evaluate rare violations, recovery, and reward hacking rather than only mean
return.

\end{v3topic}
\begin{v3topic}{Off-Policy Evidence}
Before deploying a new policy, estimate its value with importance weighting,
doubly robust estimators, or a validated simulator.
Large policy shifts create high-variance or biased estimates when action
probabilities have little overlap.
Report effective sample size, sensitivity to clipping, and where evaluation
depends on extrapolation.
\end{v3topic}


\begin{practicallab}[Watch an agent learn, and then watch it cheat]
\textbf{Goal.} Build the smallest RL system that exhibits all three classic
failure modes, so you recognise them before they cost you a quarter.

\textbf{Setup.} A $4\times4$ gridworld with a goal, one pit, and deterministic
moves. Tabular Q-learning. No libraries needed beyond an array.

\textbf{Protocol.}
\begin{enumerate}
  \item Implement Q-learning with $\gamma = 0.9$, $\alpha = 0.1$,
        $\epsilon = 0.1$. Plot return per episode until it plateaus.
  \item Print the value table and check it against
        \cref{fig:rl-value-propagation}: values should decay by roughly
        $\gamma$ per step of distance from the goal.
  \item \textbf{Sparse reward.} Move the goal to the far corner and give reward
        only there. Record how many episodes are needed now. Then set
        $\gamma = 0.99$ and repeat.
  \item \textbf{Reward hacking.} Add a small $+0.1$ for stepping on any tile
        the agent has not visited this episode. Watch what the agent does
        instead of reaching the goal.
  \item \textbf{On-policy vs off-policy.} Put a pit alongside the shortest
        path. Train Q-learning and SARSA with the same $\epsilon$. Compare the
        routes they learn and explain the difference from the update rule.
  \item Re-run step 1 with three different random seeds and plot all three.
\end{enumerate}

\textbf{Deliverable.} Four plots and one paragraph per step 3--5 explaining the
behaviour from the update equation. \textbf{The point:} step 4 takes about five
minutes to produce and is the entire reason reward design is the hardest part
of applied RL. Step 6 is why single-seed RL results should not be believed.
\end{practicallab}

\begin{checkpoint}
You are asked to use RL to set prices in a live marketplace. Give three
questions you would need answered before agreeing, drawn from this chapter:
one about the reward, one about exploration cost in a real environment, and one
about how you would evaluate the new policy before deploying it.
\end{checkpoint}

\chapterreview{Model}{Act}{Update}{Evaluate}


    \section{Further Reading}

\begin{v3topic}{Recommended Sources for Reinforcement Learning and Causality}
    \begin{itemize}
        \item \fullcite{suttonReinforcementLearningIntroduction2018} --- The definitive RL textbook: MDPs, TD learning, policy gradients, function approximation, and applications.
        \item \fullcite{hurbansGrokkingArtificialIntelligence2020} --- Accessible introduction to RL with visual examples, Q-learning, and exploration strategies.
        \item \fullcite{barberBayesianReasoningMachine2012} --- Bayesian inference, graphical models, MDPs, and value iteration with rigorous probabilistic treatment.
        \item \fullcite{downeyThinkBayes2016} --- Practical Bayesian statistics with Python: conjugate priors, MCMC, and sequential updating.
        \item \fullcite{pearlCausalityModelsReasoning2009} --- The foundational text on structural causal models, do-calculus, and causal identification.
        \item \fullcite{pearlBookWhyNew2018} --- Popular-science introduction to the ``ladder of causation'' and the do-operator.
        \item \fullcite{pearlCausalInferenceStatistics2016} --- Concise primer on causal inference: DAGs, adjustment, frontdoor/backdoor criteria.
        \item \fullcite{hernanCausalInferenceWhat2020} --- Practical causal inference for health sciences: RCTs, observational studies, time-varying treatments.
        \item \fullcite{nessCausalAI2025} --- Modern perspective on causal AI, including Granger causality, causal discovery, and missing data.
    \end{itemize}

\end{v3topic}

%